% --- Archon physics formalization source begin ---
% archon:physics
% archon:covers PhyXMiniProblems/problem_phyx_mini_0556.lean
% archon:source-report reports/phyx_mini/problem_phyx_mini_0556.source.json

\chapter{Physics problem phyx\_mini\_0556}
\label{ch:PhyXMiniProblems_problem_phyx_mini_0556}

\paragraph{Problem source.}
\#\# Physical scenario

Particle A has a rest energy of \textbackslash{}( 1192 \textbackslash{}text\{ MeV\} \textbackslash{}) and is moving through the laboratory in the positive \textbackslash{}( x \textbackslash{}) direction with a speed of \textbackslash{}( 0.45c \textbackslash{}). It decays into particle B (rest energy = \textbackslash{}( 1116 \textbackslash{}text\{ MeV\} \textbackslash{})) and a photon; particle A disappears in the decay process. Particle B moves at a speed of \textbackslash{}( 0.40c \textbackslash{}) at an angle of \textbackslash{}( 3.03\textasciicircum{}\textbackslash{}circ \textbackslash{}) with the positive \textbackslash{}( x \textbackslash{}) axis. The photon moves in a direction at an angle \textbackslash{}( \textbackslash{}theta \textbackslash{}) with the positive \textbackslash{}( x \textbackslash{}) axis as shown in figure.

\#\# Question

Find the angle \textbackslash{}( \textbackslash{}theta \textbackslash{}).

\#\# Answer choices

- A: A: \textbackslash{}( 9.8\textasciicircum{}\textbackslash{}circ \textbackslash{})

- B: B: \textbackslash{}( 10.5\textasciicircum{}\textbackslash{}circ \textbackslash{})

- C: C: \textbackslash{}( 13.3\textasciicircum{}\textbackslash{}circ \textbackslash{})

- D: D: \textbackslash{}( 12.7\textasciicircum{}\textbackslash{}circ \textbackslash{})

\#\# Figure caption (auxiliary; use the image as primary evidence)

The image is a physics diagram depicting a collision event. There are two main objects labeled A and B represented by black circles. 

- Object A is moving horizontally to the right with a velocity of \textbackslash{}(0.45c\textbackslash{}), indicated by an arrow.
- Object B is moving downward at an angle with respect to the horizontal, at \textbackslash{}(3.03\textasciicircum{}\textbackslash{}circ\textbackslash{}), with a velocity of \textbackslash{}(0.40c\textbackslash{}), also indicated by an arrow.
- A photon is emitted at an angle \textbackslash{}(\textbackslash{}theta\textbackslash{}) upward from the point of collision.
- The setup suggests a conservation of momentum scenario in a relativistic context, given the use of \textbackslash{}(c\textbackslash{}), the speed of light.

The directions and velocities are shown with arrows, and the angles are marked with lines extending from the point of collision.

\#\# Dataset metadata

Relativity; Physical Model Grounding Reasoning, Spatial Relation Reasoning

\paragraph{Recorded answer/context.}
D: D: \textbackslash{}( 12.7\textasciicircum{}\textbackslash{}circ \textbackslash{})

\paragraph{Figure/image path.}
/root/proposal\_for\_physic/science-mango/phyx\_mini\_run/phyx\_data/test\_image/556.png

\paragraph{Formalization target.}
create a compiling Lean file with sorry bodies at `PhyXMiniProblems/problem\_phyx\_mini\_0556.lean`. The Lean declarations must preserve the physical quantities, dimensions or dimensional roles, figure labels, governing-law hypotheses, and final relation expressed by this problem.
Use Mathlib/Physlib names found through LeanExplore where available. If a physics API is missing, introduce faithful local abstractions rather than scalar placeholder aliases.

\begin{theorem}[Physics formalization target]
\label{thm:physics:phyx_mini_0556:target}
\lean{PhyXMiniProblems.ProblemPhyXMini0556.problem_phyx_mini_0556}
\uses{def:physics:phyx-mini-0556:phyxminiproblems-problemphyxmini0556-relativisticdecaysetup, def:physics:phyx-mini-0556:phyxminiproblems-problemphyxmini0556-matchestwobodydecayscenario, def:physics:phyx-mini-0556:phyxminiproblems-problemphyxmini0556-matchessupplieddecayfigure, def:physics:phyx-mini-0556:phyxminiproblems-problemphyxmini0556-matchesproblemreadouts, def:physics:phyx-mini-0556:phyxminiproblems-problemphyxmini0556-hasphysicaldecayparameters, def:physics:phyx-mini-0556:phyxminiproblems-problemphyxmini0556-satisfiesrelativisticdecaylaws, def:physics:phyx-mini-0556:phyxminiproblems-problemphyxmini0556-answerchoice, def:physics:phyx-mini-0556:phyxminiproblems-problemphyxmini0556-matchesdisplayedphotonangle}
The assigned autoformalize agent should translate this physics problem into Lean declarations in the covered file, with theorem and lemma proof bodies written as `by sorry`.
\end{theorem}
\begin{proof}
This is an autoformalization task, not a proof task. Produce statements that can later be proved by the physics prover without weakening the physical model.
\end{proof}

% --- Archon declaration topology begin ---
\section{Declaration topology}
\begin{definition}[Plane Momentum]
  \label{def:physics:phyx-mini-0556:phyxminiproblems-problemphyxmini0556-planemomentum}
  \lean{PhyXMiniProblems.ProblemPhyXMini0556.PlaneMomentum}
  A unit-independent two-dimensional physical momentum in the plane of the decay diagram
\end{definition}
\begin{proof}
  This is the declared model definition of Plane Momentum.
\end{proof}
\begin{definition}[Diagram Axis]
  \label{def:physics:phyx-mini-0556:phyxminiproblems-problemphyxmini0556-diagramaxis}
  \lean{PhyXMiniProblems.ProblemPhyXMini0556.DiagramAxis}
  The two coordinate directions in the decay diagram
\end{definition}
\begin{proof}
  This is the declared model definition of Diagram Axis.
\end{proof}
\begin{definition}[to Fin]
  \label{def:physics:phyx-mini-0556:phyxminiproblems-problemphyxmini0556-diagramaxis-tofin}
  \lean{PhyXMiniProblems.ProblemPhyXMini0556.DiagramAxis.toFin}
  \uses{def:physics:phyx-mini-0556:phyxminiproblems-problemphyxmini0556-diagramaxis}
  The Fin 2 coordinate corresponding to a labelled diagram axis
\end{definition}
\begin{proof}
  This is the declared model definition of to Fin.
\end{proof}
\begin{definition}[momentum Component In SI]
  \label{def:physics:phyx-mini-0556:phyxminiproblems-problemphyxmini0556-momentumcomponentinsi}
  \lean{PhyXMiniProblems.ProblemPhyXMini0556.momentumComponentInSI}
  \uses{def:physics:phyx-mini-0556:phyxminiproblems-problemphyxmini0556-planemomentum, def:physics:phyx-mini-0556:phyxminiproblems-problemphyxmini0556-diagramaxis}
  A momentum component read in coherent SI units, hence in kg · m / s
\end{definition}
\begin{proof}
  This is the declared model definition of momentum Component In SI.
\end{proof}
\begin{definition}[momentum Magnitude In SI]
  \label{def:physics:phyx-mini-0556:phyxminiproblems-problemphyxmini0556-momentummagnitudeinsi}
  \lean{PhyXMiniProblems.ProblemPhyXMini0556.momentumMagnitudeInSI}
  \uses{def:physics:phyx-mini-0556:phyxminiproblems-problemphyxmini0556-planemomentum, def:physics:phyx-mini-0556:phyxminiproblems-problemphyxmini0556-momentumcomponentinsi}
  Euclidean magnitude of a plane momentum, read in kg · m / s
\end{definition}
\begin{proof}
  This is the declared model definition of momentum Magnitude In SI.
\end{proof}
\begin{definition}[energy In Joules]
  \label{def:physics:phyx-mini-0556:phyxminiproblems-problemphyxmini0556-energyinjoules}
  \lean{PhyXMiniProblems.ProblemPhyXMini0556.energyInJoules}
  A dimensionful energy read in coherent SI units, hence in joules
\end{definition}
\begin{proof}
  This is the declared model definition of energy In Joules.
\end{proof}
\begin{definition}[energy In Mega Electron Volts]
  \label{def:physics:phyx-mini-0556:phyxminiproblems-problemphyxmini0556-energyinmegaelectronvolts}
  \lean{PhyXMiniProblems.ProblemPhyXMini0556.energyInMegaElectronVolts}
  \uses{def:physics:phyx-mini-0556:phyxminiproblems-problemphyxmini0556-energyinjoules}
  A dimensionful energy read in megaelectronvolts
\end{definition}
\begin{proof}
  This is the declared model definition of energy In Mega Electron Volts.
\end{proof}
\begin{definition}[speed In Meters Per Second]
  \label{def:physics:phyx-mini-0556:phyxminiproblems-problemphyxmini0556-speedinmeterspersecond}
  \lean{PhyXMiniProblems.ProblemPhyXMini0556.speedInMetersPerSecond}
  A nonnegative dimensionful speed read in meters per second
\end{definition}
\begin{proof}
  This is the declared model definition of speed In Meters Per Second.
\end{proof}
\begin{definition}[vacuum Speed Of Light In Meters Per Second]
  \label{def:physics:phyx-mini-0556:phyxminiproblems-problemphyxmini0556-vacuumspeedoflightinmeterspersecond}
  \lean{PhyXMiniProblems.ProblemPhyXMini0556.vacuumSpeedOfLightInMetersPerSecond}
  Physlib's exact vacuum speed of light, read in meters per second
\end{definition}
\begin{proof}
  This is the declared model definition of vacuum Speed Of Light In Meters Per Second.
\end{proof}
\begin{definition}[speed Fraction Of Light]
  \label{def:physics:phyx-mini-0556:phyxminiproblems-problemphyxmini0556-speedfractionoflight}
  \lean{PhyXMiniProblems.ProblemPhyXMini0556.speedFractionOfLight}
  \uses{def:physics:phyx-mini-0556:phyxminiproblems-problemphyxmini0556-speedinmeterspersecond, def:physics:phyx-mini-0556:phyxminiproblems-problemphyxmini0556-vacuumspeedoflightinmeterspersecond}
  The dimensionless relativistic speed parameter β = v / c
\end{definition}
\begin{proof}
  This is the declared model definition of speed Fraction Of Light.
\end{proof}
\begin{definition}[angle In Degrees]
  \label{def:physics:phyx-mini-0556:phyxminiproblems-problemphyxmini0556-angleindegrees}
  \lean{PhyXMiniProblems.ProblemPhyXMini0556.angleInDegrees}
  Convert a radian-valued angle into its degree readout
\end{definition}
\begin{proof}
  This is the declared model definition of angle In Degrees.
\end{proof}
\begin{definition}[Rounds To Resolution]
  \label{def:physics:phyx-mini-0556:phyxminiproblems-problemphyxmini0556-roundstoresolution}
  \lean{PhyXMiniProblems.ProblemPhyXMini0556.RoundsToResolution}
  A scalar readout rounds to reported at a stated resolution
\end{definition}
\begin{proof}
  This is the declared model definition of Rounds To Resolution.
\end{proof}
\begin{definition}[Particle Label]
  \label{def:physics:phyx-mini-0556:phyxminiproblems-problemphyxmini0556-particlelabel}
  \lean{PhyXMiniProblems.ProblemPhyXMini0556.ParticleLabel}
  The three particle labels occurring in the decay
\end{definition}
\begin{proof}
  This is the declared model definition of Particle Label.
\end{proof}
\begin{definition}[Massive Particle]
  \label{def:physics:phyx-mini-0556:phyxminiproblems-problemphyxmini0556-massiveparticle}
  \lean{PhyXMiniProblems.ProblemPhyXMini0556.MassiveParticle}
  The two massive particles to which E = γ E₀ applies
\end{definition}
\begin{proof}
  This is the declared model definition of Massive Particle.
\end{proof}
\begin{definition}[to Particle Label]
  \label{def:physics:phyx-mini-0556:phyxminiproblems-problemphyxmini0556-massiveparticle-toparticlelabel}
  \lean{PhyXMiniProblems.ProblemPhyXMini0556.MassiveParticle.toParticleLabel}
  \uses{def:physics:phyx-mini-0556:phyxminiproblems-problemphyxmini0556-particlelabel, def:physics:phyx-mini-0556:phyxminiproblems-problemphyxmini0556-massiveparticle}
  Regard a massive-particle label as a label in the full decay event
\end{definition}
\begin{proof}
  This is the declared model definition of to Particle Label.
\end{proof}
\begin{definition}[Decay Stage]
  \label{def:physics:phyx-mini-0556:phyxminiproblems-problemphyxmini0556-decaystage}
  \lean{PhyXMiniProblems.ProblemPhyXMini0556.DecayStage}
  State immediately before or immediately after the decay
\end{definition}
\begin{proof}
  This is the declared model definition of Decay Stage.
\end{proof}
\begin{definition}[Path Side]
  \label{def:physics:phyx-mini-0556:phyxminiproblems-problemphyxmini0556-pathside}
  \lean{PhyXMiniProblems.ProblemPhyXMini0556.PathSide}
  Which side of the horizontal axis contains a depicted particle path
\end{definition}
\begin{proof}
  This is the declared model definition of Path Side.
\end{proof}
\begin{definition}[Horizontal Direction]
  \label{def:physics:phyx-mini-0556:phyxminiproblems-problemphyxmini0556-horizontaldirection}
  \lean{PhyXMiniProblems.ProblemPhyXMini0556.HorizontalDirection}
  Direction of the positive horizontal reference axis
\end{definition}
\begin{proof}
  This is the declared model definition of Horizontal Direction.
\end{proof}
\begin{definition}[Relativistic Decay Figure]
  \label{def:physics:phyx-mini-0556:phyxminiproblems-problemphyxmini0556-relativisticdecayfigure}
  \lean{PhyXMiniProblems.ProblemPhyXMini0556.RelativisticDecayFigure}
  \uses{def:physics:phyx-mini-0556:phyxminiproblems-problemphyxmini0556-particlelabel, def:physics:phyx-mini-0556:phyxminiproblems-problemphyxmini0556-pathside, def:physics:phyx-mini-0556:phyxminiproblems-problemphyxmini0556-horizontaldirection}
  Literal labels and geometric annotations visible in the supplied image
\end{definition}
\begin{proof}
  This is the declared model definition of Relativistic Decay Figure.
\end{proof}
\begin{definition}[Relativistic Decay Setup]
  \label{def:physics:phyx-mini-0556:phyxminiproblems-problemphyxmini0556-relativisticdecaysetup}
  \lean{PhyXMiniProblems.ProblemPhyXMini0556.RelativisticDecaySetup}
  \uses{def:physics:phyx-mini-0556:phyxminiproblems-problemphyxmini0556-planemomentum, def:physics:phyx-mini-0556:phyxminiproblems-problemphyxmini0556-particlelabel, def:physics:phyx-mini-0556:phyxminiproblems-problemphyxmini0556-decaystage, def:physics:phyx-mini-0556:phyxminiproblems-problemphyxmini0556-relativisticdecayfigure}
  The physical quantities in the event. Rest and total energies are distinct fields. The photon angle is likewise an independent observable rather than a definition involving an answer choice
\end{definition}
\begin{proof}
  This is the declared model definition of Relativistic Decay Setup.
\end{proof}
\begin{definition}[massive Rest Energy]
  \label{def:physics:phyx-mini-0556:phyxminiproblems-problemphyxmini0556-massiverestenergy}
  \lean{PhyXMiniProblems.ProblemPhyXMini0556.massiveRestEnergy}
  \uses{def:physics:phyx-mini-0556:phyxminiproblems-problemphyxmini0556-massiveparticle, def:physics:phyx-mini-0556:phyxminiproblems-problemphyxmini0556-relativisticdecaysetup}
  Rest energy of one of the two massive particles
\end{definition}
\begin{proof}
  This is the declared model definition of massive Rest Energy.
\end{proof}
\begin{definition}[massive Speed]
  \label{def:physics:phyx-mini-0556:phyxminiproblems-problemphyxmini0556-massivespeed}
  \lean{PhyXMiniProblems.ProblemPhyXMini0556.massiveSpeed}
  \uses{def:physics:phyx-mini-0556:phyxminiproblems-problemphyxmini0556-massiveparticle, def:physics:phyx-mini-0556:phyxminiproblems-problemphyxmini0556-relativisticdecaysetup}
  Laboratory-frame speed of one of the two massive particles
\end{definition}
\begin{proof}
  This is the declared model definition of massive Speed.
\end{proof}
\begin{definition}[massive Total Energy]
  \label{def:physics:phyx-mini-0556:phyxminiproblems-problemphyxmini0556-massivetotalenergy}
  \lean{PhyXMiniProblems.ProblemPhyXMini0556.massiveTotalEnergy}
  \uses{def:physics:phyx-mini-0556:phyxminiproblems-problemphyxmini0556-massiveparticle, def:physics:phyx-mini-0556:phyxminiproblems-problemphyxmini0556-relativisticdecaysetup}
  Laboratory-frame total energy of one of the two massive particles
\end{definition}
\begin{proof}
  This is the declared model definition of massive Total Energy.
\end{proof}
\begin{definition}[massive Momentum]
  \label{def:physics:phyx-mini-0556:phyxminiproblems-problemphyxmini0556-massivemomentum}
  \lean{PhyXMiniProblems.ProblemPhyXMini0556.massiveMomentum}
  \uses{def:physics:phyx-mini-0556:phyxminiproblems-problemphyxmini0556-planemomentum, def:physics:phyx-mini-0556:phyxminiproblems-problemphyxmini0556-massiveparticle, def:physics:phyx-mini-0556:phyxminiproblems-problemphyxmini0556-relativisticdecaysetup}
  Laboratory-frame plane momentum of one of the two massive particles
\end{definition}
\begin{proof}
  This is the declared model definition of massive Momentum.
\end{proof}
\begin{definition}[Matches Two Body Decay Scenario]
  \label{def:physics:phyx-mini-0556:phyxminiproblems-problemphyxmini0556-matchestwobodydecayscenario}
  \lean{PhyXMiniProblems.ProblemPhyXMini0556.MatchesTwoBodyDecayScenario}
  \uses{def:physics:phyx-mini-0556:phyxminiproblems-problemphyxmini0556-relativisticdecaysetup}
  Particle A exists before the decay and disappears; particle B and the photon are the two products present afterward
\end{definition}
\begin{proof}
  This is the declared model definition of Matches Two Body Decay Scenario.
\end{proof}
\begin{definition}[Matches Supplied Decay Figure]
  \label{def:physics:phyx-mini-0556:phyxminiproblems-problemphyxmini0556-matchessupplieddecayfigure}
  \lean{PhyXMiniProblems.ProblemPhyXMini0556.MatchesSuppliedDecayFigure}
  \uses{def:physics:phyx-mini-0556:phyxminiproblems-problemphyxmini0556-momentumcomponentinsi, def:physics:phyx-mini-0556:phyxminiproblems-problemphyxmini0556-momentummagnitudeinsi, def:physics:phyx-mini-0556:phyxminiproblems-problemphyxmini0556-relativisticdecaysetup}
  Primary-image evidence. Besides the literal labels and branch placement, the component equations state what it means for the displayed angles to be measured above or below the positive x axis
\end{definition}
\begin{proof}
  This is the declared model definition of Matches Supplied Decay Figure.
\end{proof}
\begin{definition}[Matches Problem Readouts]
  \label{def:physics:phyx-mini-0556:phyxminiproblems-problemphyxmini0556-matchesproblemreadouts}
  \lean{PhyXMiniProblems.ProblemPhyXMini0556.MatchesProblemReadouts}
  \uses{def:physics:phyx-mini-0556:phyxminiproblems-problemphyxmini0556-energyinmegaelectronvolts, def:physics:phyx-mini-0556:phyxminiproblems-problemphyxmini0556-speedfractionoflight, def:physics:phyx-mini-0556:phyxminiproblems-problemphyxmini0556-angleindegrees, def:physics:phyx-mini-0556:phyxminiproblems-problemphyxmini0556-roundstoresolution, def:physics:phyx-mini-0556:phyxminiproblems-problemphyxmini0556-relativisticdecaysetup}
  Numerical problem data. The rest energies and speed ratios are the stated idealized values. The displayed 3.03° direction is a two-decimal readout of the physical B direction; this avoids treating a rounded annotation as an additional exact conservation equation
\end{definition}
\begin{proof}
  This is the declared model definition of Matches Problem Readouts.
\end{proof}
\begin{definition}[Has Physical Decay Parameters]
  \label{def:physics:phyx-mini-0556:phyxminiproblems-problemphyxmini0556-hasphysicaldecayparameters}
  \lean{PhyXMiniProblems.ProblemPhyXMini0556.HasPhysicalDecayParameters}
  \uses{def:physics:phyx-mini-0556:phyxminiproblems-problemphyxmini0556-momentummagnitudeinsi, def:physics:phyx-mini-0556:phyxminiproblems-problemphyxmini0556-energyinjoules, def:physics:phyx-mini-0556:phyxminiproblems-problemphyxmini0556-vacuumspeedoflightinmeterspersecond, def:physics:phyx-mini-0556:phyxminiproblems-problemphyxmini0556-speedfractionoflight, def:physics:phyx-mini-0556:phyxminiproblems-problemphyxmini0556-massiveparticle, def:physics:phyx-mini-0556:phyxminiproblems-problemphyxmini0556-relativisticdecaysetup, def:physics:phyx-mini-0556:phyxminiproblems-problemphyxmini0556-massivespeed, def:physics:phyx-mini-0556:phyxminiproblems-problemphyxmini0556-massivemomentum}
  Positivity, subluminality, and the acute-angle branches shown by the figure. These select the physical trigonometric branches without fixing the unknown photon angle numerically
\end{definition}
\begin{proof}
  This is the declared model definition of Has Physical Decay Parameters.
\end{proof}
\begin{definition}[Satisfies Relativistic Decay Laws]
  \label{def:physics:phyx-mini-0556:phyxminiproblems-problemphyxmini0556-satisfiesrelativisticdecaylaws}
  \lean{PhyXMiniProblems.ProblemPhyXMini0556.SatisfiesRelativisticDecayLaws}
  \uses{def:physics:phyx-mini-0556:phyxminiproblems-problemphyxmini0556-diagramaxis, def:physics:phyx-mini-0556:phyxminiproblems-problemphyxmini0556-momentumcomponentinsi, def:physics:phyx-mini-0556:phyxminiproblems-problemphyxmini0556-momentummagnitudeinsi, def:physics:phyx-mini-0556:phyxminiproblems-problemphyxmini0556-energyinjoules, def:physics:phyx-mini-0556:phyxminiproblems-problemphyxmini0556-vacuumspeedoflightinmeterspersecond, def:physics:phyx-mini-0556:phyxminiproblems-problemphyxmini0556-speedfractionoflight, def:physics:phyx-mini-0556:phyxminiproblems-problemphyxmini0556-massiveparticle, def:physics:phyx-mini-0556:phyxminiproblems-problemphyxmini0556-relativisticdecaysetup, def:physics:phyx-mini-0556:phyxminiproblems-problemphyxmini0556-massiverestenergy, def:physics:phyx-mini-0556:phyxminiproblems-problemphyxmini0556-massivespeed, def:physics:phyx-mini-0556:phyxminiproblems-problemphyxmini0556-massivetotalenergy, def:physics:phyx-mini-0556:phyxminiproblems-problemphyxmini0556-massivemomentum}
  For each massive particle, E = γ(β) E₀ and p c = γ(β) β E₀. The photon obeys E = p c; total energy and both components of momentum are conserved at the decay vertex. These laws are generic in the independent physical fields and contain no solved angle or answer choice
\end{definition}
\begin{proof}
  This is the declared model definition of Satisfies Relativistic Decay Laws.
\end{proof}
\begin{lemma}[photon Angle tan from momentum Conservation]
  \label{lem:physics:phyx-mini-0556:phyxminiproblems-problemphyxmini0556-photonangle-tan-from-momentumconservation}
  \lean{PhyXMiniProblems.ProblemPhyXMini0556.photonAngle_tan_from_momentumConservation}
  \uses{def:physics:phyx-mini-0556:phyxminiproblems-problemphyxmini0556-momentummagnitudeinsi, def:physics:phyx-mini-0556:phyxminiproblems-problemphyxmini0556-relativisticdecaysetup, def:physics:phyx-mini-0556:phyxminiproblems-problemphyxmini0556-matchessupplieddecayfigure, def:physics:phyx-mini-0556:phyxminiproblems-problemphyxmini0556-hasphysicaldecayparameters, def:physics:phyx-mini-0556:phyxminiproblems-problemphyxmini0556-satisfiesrelativisticdecaylaws}
  Momentum conservation determines the tangent of the photon angle from the independently represented A and B momenta and the displayed branch geometry
\end{lemma}
\begin{proof}
  The conclusion follows from the governing relations and figure readouts named in the statement of photon Angle tan from momentum Conservation.
\end{proof}
\begin{definition}[Answer Choice]
  \label{def:physics:phyx-mini-0556:phyxminiproblems-problemphyxmini0556-answerchoice}
  \lean{PhyXMiniProblems.ProblemPhyXMini0556.AnswerChoice}
  Labels of the four angle choices printed in the problem
\end{definition}
\begin{proof}
  This is the declared model definition of Answer Choice.
\end{proof}
\begin{definition}[displayed Photon Angle Degrees]
  \label{def:physics:phyx-mini-0556:phyxminiproblems-problemphyxmini0556-displayedphotonangledegrees}
  \lean{PhyXMiniProblems.ProblemPhyXMini0556.displayedPhotonAngleDegrees}
  \uses{def:physics:phyx-mini-0556:phyxminiproblems-problemphyxmini0556-answerchoice}
  Degree-valued angle printed beside each answer label
\end{definition}
\begin{proof}
  This is the declared model definition of displayed Photon Angle Degrees.
\end{proof}
\begin{definition}[recorded Dataset Answer]
  \label{def:physics:phyx-mini-0556:phyxminiproblems-problemphyxmini0556-recordeddatasetanswer}
  \lean{PhyXMiniProblems.ProblemPhyXMini0556.recordedDatasetAnswer}
  \uses{def:physics:phyx-mini-0556:phyxminiproblems-problemphyxmini0556-answerchoice}
  The answer label recorded in the source dataset, retained only as metadata and not used as a theorem premise
\end{definition}
\begin{proof}
  This is the declared model definition of recorded Dataset Answer.
\end{proof}
\begin{definition}[Matches Displayed Photon Angle]
  \label{def:physics:phyx-mini-0556:phyxminiproblems-problemphyxmini0556-matchesdisplayedphotonangle}
  \lean{PhyXMiniProblems.ProblemPhyXMini0556.MatchesDisplayedPhotonAngle}
  \uses{def:physics:phyx-mini-0556:phyxminiproblems-problemphyxmini0556-angleindegrees, def:physics:phyx-mini-0556:phyxminiproblems-problemphyxmini0556-roundstoresolution, def:physics:phyx-mini-0556:phyxminiproblems-problemphyxmini0556-relativisticdecaysetup, def:physics:phyx-mini-0556:phyxminiproblems-problemphyxmini0556-answerchoice, def:physics:phyx-mini-0556:phyxminiproblems-problemphyxmini0556-displayedphotonangledegrees}
  The physical photon angle rounds to the selected one-decimal-degree answer
\end{definition}
\begin{proof}
  This is the declared model definition of Matches Displayed Photon Angle.
\end{proof}
% --- Archon declaration topology end ---
% --- Archon physics formalization source end ---
